% =============================================================================
% P4 --- Discusion reescrita para integrar las figuras nuevas
% (stage2_morfologia_mavic_hunter y snr_ejemplos) y unificar las dos partes
% bajo la hipotesis morfologica. Sustituye el \subsection{Discusión} de P4.
% Nota: en la seccion de resultados, corregir el doble punto final de
% "...no un rediseño del sistema.." -> "...no un rediseño del sistema."
% =============================================================================
\subsection{Discusión}
\label{sec:p4_discusion}

\paragraph{Identificación del modelo de dron.}
La primera parte de P4 confirma que el sistema no solo detecta la presencia de un dron, sino que identifica su modelo concreto. El DJI Mini~3 ---la plataforma objetivo de este trabajo--- se clasifica correctamente en prácticamente la totalidad de las muestras (F1\,=\,\num{0,992}), y la clase \emph{interference} obtiene precisión \num{1,000}: ninguna señal real de interferencia se confunde de forma sistemática con un dron, en línea con la ausencia de confusión cruzada observada en P2 y P3. Los \num{25}~errores sobre \num{300}~muestras de interferencia se reparten entre varias clases de dron sin concentrarse en ninguna, lo que descarta un sesgo sistemático.

El único fallo relevante ---el colapso de \emph{mavic\_novideo} sobre \emph{hunter}--- no contradice la hipótesis central del trabajo, sino que la refuerza. Como muestra la \cref{fig:stage2_morfologia}, cuando el Mavic deja de transmitir vídeo su emisión se reduce a la señal de control FHSS, morfológicamente indistinguible de la del Hunter. Que el clasificador agrupe ambas clases es precisamente lo que cabe esperar de una red que decide por la \emph{forma} de la señal y no por una etiqueta o por el contexto de captura: ante dos morfologías idénticas, emite la misma predicción. El error no es, por tanto, un defecto del modelo, sino una limitación de los datos ---la ausencia, en ese modo de emisión, de cualquier rasgo espectral que separe las dos plataformas---, resoluble con más capturas de la clase en condiciones variadas y no con un rediseño del sistema.

\paragraph{Cómo se inyecta el ruido.}
Conviene recordar cómo se construyen ambos ruidos, descritos en detalle en la \cref{sec:implementacion_snr}, porque su forma de generación es la clave de la comparación. Los dos se inyectan en el dominio IQ complejo de la señal y se calibran a la \emph{misma potencia total} para cada SNR objetivo, de modo que la única variable que cambia entre ellos es cómo se reparte esa energía en el plano tiempo-frecuencia. El \textbf{ruido blanco AWGN} es ruido gaussiano complejo, $n[k]\sim\mathcal{CN}(0,\sigma^2)$, sumado muestra a muestra a la señal con la varianza $\sigma^2$ ajustada al SNR objetivo; al ser blanco, su energía se distribuye de forma uniforme por toda la banda y toda la ventana, elevando el fondo del espectrograma por igual en todo el plano. El \textbf{ruido FHSS sintético} emula una interferencia tipo Bluetooth Classic: genera unas \num{160}~ráfagas por ventana de \SI{0.1}{\second} (\num{1600}~saltos/s), cada una de \SI{366}{\micro\second} de duración y aproximadamente \SI{1}{\mega\hertz} de ancho, colocada en un instante y una frecuencia aleatorios dentro de la banda capturada y suavizada con una ventana de Hann. Como el conjunto se calibra a la misma potencia total que el AWGN, esa energía idéntica queda concentrada en rectángulos estrechos y dispersos ---morfológicamente parecidos a las propias ráfagas del dron--- que ocupan solo en torno al \SI{1,5}{\percent} del plano tiempo-frecuencia.

\paragraph{Robustez frente al ruido.}
La segunda parte de P4 somete al clasificador a dos modelos de ruido de la misma potencia total pero distinta distribución en el plano tiempo-frecuencia. El contraste es nítido tanto a nivel visual como cuantitativo. Visualmente (\cref{fig:snr_ejemplos}), el AWGN eleva el fondo de forma uniforme hasta enmascarar por completo las ráfagas a $-20$~dB, mientras que el FHSS sintético deja intacto el grueso del espectrograma y mantiene visibles las ráfagas de OcuSync incluso a SNR muy bajo. Cuantitativamente, esa diferencia se traduce en el codo abrupto del AWGN entre $+5$ y $-5$~dB ---una caída de \num{0,61}~puntos de exactitud--- frente a la degradación suave del FHSS, que conserva una exactitud superior a \num{0,85} hasta $-5$~dB.

Esta asimetría solo admite una explicación: la red discrimina por \emph{dónde} está la energía, no por \emph{cuánta} hay. El AWGN reparte su potencia por los \SI{40}{\mega\hertz} del espectrograma y, al superar el nivel de la señal, borra la firma de las ráfagas; el FHSS sintético concentra la misma potencia en aproximadamente el \SI{1,5}{\percent} del plano, de modo que el \SI{98,5}{\percent} restante ---y con él la morfología del dron--- permanece intacto. Si el clasificador se apoyase en el nivel medio de potencia, ambos ruidos lo degradarían por igual, dado que inyectan la misma energía total; que el AWGN sea mucho más destructivo demuestra lo contrario.

\paragraph{Síntesis: la red discrimina por morfología.}
Las dos partes de P4 convergen en la misma conclusión y cierran la hipótesis central del trabajo. La confusión \emph{mavic\_novideo}--\emph{hunter} muestra que la red agrupa señales por su morfología espectro-temporal; el barrido de SNR muestra que su decisión sobrevive mientras esa morfología sea visible y se desploma solo cuando el ruido la borra. Ambas observaciones son coherentes con los mapas de GradCAM del experimento P1 y de EigenCAM del experimento P2, donde las activaciones de la red se localizaban sobre las estructuras de ráfaga y no sobre el fondo. El clasificador no aprende la intensidad media de la captura ni un atajo contextual, sino la posición de la energía en el plano tiempo-frecuencia: la firma morfológica del dron. Este resultado, sólido para el DJI Mini~3 sobre el conjunto de prueba, sustenta la tesis del trabajo, con la cautela de que su validación en condiciones de despliegue más diversas queda planteada como línea futura.
